%-------------------------
% Resume in Latex
% Author : Sourabh Bajaj
% License : MIT
%------------------------

\documentclass[letterpaper,11pt]{article}

\usepackage{latexsym}
\usepackage[empty]{fullpage}
\usepackage{titlesec}
\usepackage{marvosym}
\usepackage[usenames,dvipsnames]{color}
\usepackage{verbatim}
\usepackage{enumitem}
\usepackage[hidelinks]{hyperref}
\usepackage{fancyhdr}
\usepackage[english]{babel}
\usepackage{makecell}

\pagestyle{fancy}
\fancyhf{} % clear all header and footer fields
\fancyfoot{}
\renewcommand{\headrulewidth}{0pt}
\renewcommand{\footrulewidth}{0pt}

% Adjust margins
\addtolength{\oddsidemargin}{-0.5in}
\addtolength{\evensidemargin}{-0.5in}
\addtolength{\textwidth}{1in}
\addtolength{\topmargin}{-.5in}
\addtolength{\textheight}{1.0in}

\urlstyle{same}

\raggedbottom
\raggedright
\setlength{\tabcolsep}{0in}

% Sections formatting
\titleformat{\section}{
  \vspace{-4pt}\scshape\raggedright\large
}{}{0em}{}[\color{black}\titlerule \vspace{-5pt}]

%-------------------------
% Custom commands
\newcommand{\resumeItem}[2]{
  \item\small{
    \textbf{#1}{ #2 \vspace{-2pt}}
  }
}

\newcommand{\resumeSubheading}[4]{
  \vspace{-1pt}\item
    \begin{tabular*}{0.97\textwidth}[t]{l@{\extracolsep{\fill}}r}
      \textbf{#1} & #2 \\
      \textit{\small#3} & \textit{\small #4} \\
    \end{tabular*}\vspace{-5pt}
}

\newcommand{\resumeSubSubheading}[2]{
  \vspace{-1pt}\item
    \begin{tabular*}{0.97\textwidth}[t]{l@{\extracolsep{\fill}}r}
      \textbf{#1} & #2 \\
      %\textit{\small#3} & \textit{\small #4} \\
    \end{tabular*}\vspace{-5pt}
}

\newcommand{\resumeSubItem}[2]{\resumeItem{#1}{#2}\vspace{-4pt}}

\renewcommand{\labelitemii}{$\circ$}

\newcommand{\resumeSubHeadingListStart}{\begin{itemize}[leftmargin=*]}
\newcommand{\resumeSubHeadingListEnd}{\end{itemize}}
\newcommand{\resumeItemListStart}{\begin{itemize}}
\newcommand{\resumeItemListEnd}{\end{itemize}\vspace{-5pt}}

%-------------------------------------------
%%%%%%  CV STARTS HERE  %%%%%%%%%%%%%%%%%%%%%%%%%%%%


\begin{document}

%----------HEADING-----------------
\begin{tabular*}{\textwidth}{l@{\extracolsep{\fill}}r}
  \textbf{\href{https://ukulililixl.github.io/}{\Large Li Xiaolu}} & Email : \href{mailto:ukulililixl@gmail.com}{ukulililixl@gmail.com}\\
  \href{https://ukulililixl.github.io/}{https://ukulililixl.github.io} & Location: Hong Kong \\
\end{tabular*}


%-----------EDUCATION-----------------
\section{Education}
  \resumeSubHeadingListStart
    \resumeSubheading
      {The Chinese University of Hong Kong}{Hong Kong}
      {PhD of Computer Science and Engineering;  }{Aug. 2016 -- Now}
    \resumeSubheading
      {University of Science and Technology of China}{Anhui, China}
      {Bachelor of Computer Science and Technology;  GPA: 3.71 (88.03/100.00)}{Aug. 2012 -- June. 2016}
  \resumeSubHeadingListEnd

%-----------Publication---------------
\section{Publication}
  \resumeSubHeadingListStart
    \resumeSubheading
      {\makecell[l]{Repair Pipelining for Erasure-Coded Storage: Algorithms and Evaluation.}}{}
      {\makecell[l]{{\bf Xiaolu Li}, Zuoru Yang, Jinhong Li, Runhui Li, Patrick P. C. Lee, Qun Huang and Yuchong Hu.}}{In submission.}
    \resumeSubheading
      {\makecell[l]{Fast Predictive Repair in Erasure-Coded Storage}}{}
      {\makecell[l]{Zhirong Shen, {\bf Xiaolu Li} and Patrick P. C. Lee. \\IEEE/IFIP DSN 2019. (AR: 54/252 = 21.4\%)}}{}
    \resumeSubheading
      {\makecell[l]{OpenEC: Toward Unified and Configurable Erasure Coding Management in \\Distributed Storage Systems}}{}
      {\makecell[l]{{\bf Xiaolu Li}, Runhui Li, Patrick P. C. Lee, and Yuchong Hu. \\USENIX FAST 2019. (AR: 26/145 = 17.9\%)}}{}
    \resumeSubheading
      {Optimal Repair Layering for Erasure-Coded Data Centers: From Theory to Practice}{}
      {\makecell[l]{Yuchong Hu, {\bf Xiaolu Li}, Mi Zhang, Patrick P. C. Lee, Xiaoyang Zhang, Pan Zhou, and Dan Feng. \\ACM Transaction on Storage  (TOS), 13(4), pp. 33:1-33:24,
      November 2017}}{}
    \resumeSubheading
      {Repair Pipelining for Erasure-Coded Storage}{}
      {\makecell[l]{Runhui Li, {\bf Xiaolu Li}, Patrick P.C. Lee and Qun Huang. \\USENIX ATC 2017. (AR: 60/283 = 21.2\%)}}{}


  \resumeSubHeadingListEnd

%-----------Project-------------------
\section{Selected Projects}
  \resumeSubHeadingListStart
    \resumeSubheading
      {OpenEC}{Aug. 2017 - Sept. 2018}
      {Source: http://adslab.cse.cuhk.edu.hk/software/openec/}{}
      \resumeItemListStart
        \resumeItem{}{OpenEC is a unified and configurable erasure coding management system that easily runs atop multiple distributed storage systems to manage erasure-coded data. It presents simple programming model for erasure coding designers to design new erasure coding solutions, including new erasure codes and repair algorithms. It also provides general interfaces for integration into storage systems.}
	\resumeItem{}{I build OpenEC prototype and validate it atop Hadoop 3.0 HDFS, HDFS-RAID and QuantcastFS.}
        %\resumeItem{}{We build a recovery prototype called DoubleR atop HDFS and apply a new erasure code called Double Regenerating Code to improve the repair performance in hierarchical datacenters. We apply the layering technique to aggregate traffic in the same region before data are transfered cross region to reduce cross-region traffic.}
	%\resumeItem{}{We implement many state-of-the-art erasure codes in HDFS and shows that their repair performance is improved using our recovery prototype.}
      \resumeItemListEnd
    \resumeSubheading
      {DoubleR}{Feb. 2016 - March. 2017}
      {Source: http://adslab.cse.cuhk.edu.hk/software/doubler/}{}
      \resumeItemListStart
        \resumeItem{}{DoubleR is a recovery prototype to improve repair performance in hierarchical datacenters.}
	\resumeItem{}{I build the prototype and improve repair performance for state-of-the-art erasure codes in HDFS-RAID.}
        %\resumeItem{}{We build a recovery prototype called DoubleR atop HDFS and apply a new erasure code called Double Regenerating Code to improve the repair performance in hierarchical datacenters. We apply the layering technique to aggregate traffic in the same region before data are transfered cross region to reduce cross-region traffic.}
	%\resumeItem{}{We implement many state-of-the-art erasure codes in HDFS and shows that their repair performance is improved using our recovery prototype.}
      \resumeItemListEnd
    \resumeSubheading
      {ECPipe}{Sept. 2016 - Feb. 2017}
      {Source: http://adslab.cse.cuhk.edu.hk/software/ecpipe/} {}
      \resumeItemListStart
        \resumeItem{}{ECPipe enables constant time repair operation in erasure-coded distributed storage system.}
        \resumeItem{}{I integrate our ECPipe into HDFS-RAID as well as QuantcastFS and improve the repair performance.}
        %\resumeItem{}{We build a prototype called ECPipe which borrows the idea of cache-and-forward from network into distributed storage system for the purpose to improve the repair performance. We improve the repair performance from $O(k)$ to $O(1)$ in homogeneous network settings.}
	%\resumeItem{}{We integrate our prototype ECPipe into large scale distributed storage systems including HDFS and QuantcastFS. The repair performance in both of the two systems are improved using our recovery prototype.}
      \resumeItemListEnd

  \resumeSubHeadingListEnd

%-----------Internship----------------
\section{Internship}
\resumeSubHeadingListStart
    \resumeSubheading
      {Serverless Group, Aliware, Alibaba}{July. 2019 - Aug. 2019}
      {Role: Software developer}{}
      \resumeItemListStart
        \resumeItem{}{I design a simulator, {\bf SimServerless}, special for the serverless platform based on an
            open-source cloud simulation tool (CloudSim). SimServerless provides interface for custom traffic prediction 
            algorithms as well as auto-scaling algorithms.}
        \resumeItem{}{I analysis two workloads and find it possible to use statistical methods to predict workload with
        periodicity.}
        \resumeItem{}{I design a prediction-based, two-level scaling algorithm to schedule the containers from three
        status, including stop, standby and running. Simulation results show that this algorithm improves the resource
        utilization.}
      \resumeItemListEnd
  \resumeSubHeadingListEnd
%-----------Teaching------------------
 \section{Teaching}
  \resumeSubHeadingListStart
    \resumeSubItem{Operating System:}{Multi-threading.}
    \resumeSubItem{Introduction to Cloud Computing and Storage:}{PageRank algorithm. Data deduplication.}
    \resumeSubItem{Data Communication and Computer Networks:}{Go-back-N protocol.}
  \resumeSubHeadingListEnd

%-----------Skill----------------------
\section{Skills}
  \resumeSubHeadingListStart
    \resumeSubItem{Programming Language:}{C/C++. Java. Python. R.}
    \resumeSubItem{Distributed Storage Systems:}{HDFS. QuantcastFS. OpenStack Swift. Crail. }
    \resumeSubItem{In-Memory Systems:}{Redis. Alluxio.}
    \resumeSubItem{Others:}{OpenWhisk. A little Kubernetes.}
  \resumeSubHeadingListEnd

\section{Selected Award}
  \resumeSubHeadingListStart
    \resumeSubItem{Student Cluster Competition:}{Ranking $3^{rd}$ (team work),  ACM/IEEE SC14, New Orleans}
  \resumeSubHeadingListEnd

\section{Interest}
  \resumeSubHeadingListStart
  \resumeSubItem{}{I'm interested in distributed systems, especially the distributed storage systems.}
  \resumeSubItem{}{I'm curious about the design of cloud-native storage system in the future for various serverless applications.}
  \resumeSubItem{}{I feel excited to define and solve system problems.}
  \resumeSubItem{}{I'm curious to learn new techniques.}
  \resumeSubHeadingListEnd
    
%-----------EXPERIENCE-----------------
% \section{Experience}
%   \resumeSubHeadingListStart
% 
%     \resumeSubheading
%       {Google}{Mountain View, CA}
%       {Software Engineer}{Oct 2016 - Present}
%       \resumeItemListStart
%         \resumeItem{Tensorflow}
%           {TensorFlow is an open source software library for numerical computation using data flow graphs; primarily used for training deep learning models. Worked on APIs and performance for training models on Tensor Processing Units (TPU).}
%         \resumeItem{Apache Beam}
%           {Apache Beam is a unified model for defining both batch and streaming data-parallel processing pipelines, as well as a set of language-specific SDKs for constructing pipelines and runners.}
%       \resumeItemListEnd
% 
%     \resumeSubheading
%       {Coursera}{Mountain View, CA}
%       {Senior Software Engineer}{Jan 2014 - Oct 2016}
%       \resumeItemListStart
%         \resumeItem{Notifications}
%           {Service for sending email, push and in-app notifications. Involved in features such as delivery time optimization, tracking, queuing and A/B testing. Built an internal app to run batch campaigns for marketing etc.}
%         \resumeItem{Nostos}
%           {Bulk data processing and injection service from Hadoop to Cassandra and provides a thin REST layer on top for serving offline computed data online.}
%         \resumeItem{Workflows}
%           {Dataduct an open source workflow framework to create and manage data pipelines leveraging reusables patterns to expedite developer productivity.}
%         \resumeItem{Data Collection}
%           {Designed the internal survey and crowd sourcing platform which allowed for creating various tasks for crowd sourcing or embedding surveys across the Coursera platform.}
%         \resumeItem{Dev Environment}
%           {Analytics environment based on docker and AWS, standardized the python and R dependencies. Wrote the core libraries that are shared by all data scientists.}
%         \resumeItem{Data Warehousing}
%           {Setup, schema design and management of Amazon Redshift. Built an internal app for access to the data using a web interface. Dataduct integration for daily ETL injection into Redshift.}
%         \resumeItem{Recommendations}
%           {Core service for all recommendation systems at Coursera, currently used on the homepage and throughout the content discovery process. Worked on both offline training and online serving.}
%         \resumeItem{Content Discovery}
%           {Improved content discovery by building a new onboarding experience on coursera. Using this to personalize the search and browse experience. Also worked on ranking and indexing improvements.}
%         \resumeItem{Course Dashboards}
%           {Instructor dashboards and learner surveying tools, which helped instructors run their class better by providing data on Assignments and Learner Activity.}
%       \resumeItemListEnd
% 
%     \resumeSubheading
%       {Lucena Research}{Atlanta, GA}
%       {Data Scientist}{Summer 2012 and 2013}
%       \resumeItemListStart
%         \resumeItem{Portfolio Management}
%           {Created models for portfolio hedging,  portfolio optimization and price forecasting. Also creating a strategy backtesting engine used for simulating and backtesting strategies.}
%         \resumeItem{QuantDesk}
%           {Python backend for a web application used by hedge fund managers for portfolio management.}
%       \resumeItemListEnd
% 
%     \resumeSubheading
%       {Georgia Institute of Technology}{Atlanta, GA}
%       {Research and Teaching Assistant}{Jan 2012 - Dec 2013}
%       \resumeItemListStart
%         \resumeItem{Research Assistant - Machine Learning}
%           {Research on machine learning for portfolio hedging and replication algorithms. Modeling low-risk \& continuous-return strategies. Developed the python library QSTK.}
%         \resumeItem{Teaching Assistant - Computational Investing}
%           {The online course on Coursera, had more than 100,000 students enrolled. It was featured on the 11 Alive News and the Atlanta Journal Constitution. Involved in creating assignment, exams and conducting recitation sessions. Also taught the on-campus version of the course.}
%       \resumeItemListEnd
% 
%   \resumeSubHeadingListEnd
% 
% 
% %-----------PROJECTS-----------------
% \section{Projects}
%   \resumeSubHeadingListStart
%     \resumeSubItem{QuantSoftware Toolkit}
%       {Open source python library for financial data analysis and machine learning for finance.}
%     \resumeSubItem{Github Visualization}
%       {Data Visualization of Git Log data using D3 to analyze project trends over time.}
%     \resumeSubItem{Recommendation System}
%       {Music and Movie recommender systems using collaborative filtering on public datasets.}
%     \resumeSubItem{Mac Setup}
%       {Book that gives step by step instructions on setting up developer environment on Mac OS.}
%   \resumeSubHeadingListEnd
% 
% %
% %--------PROGRAMMING SKILLS------------
% %\section{Programming Skills}
% %  \resumeSubHeadingListStart
% %    \item{
% %      \textbf{Languages}{: Scala, Python, Javascript, C++, SQL, Java}
% %      \hfill
% %      \textbf{Technologies}{: AWS, Play, React, Kafka, GCE}
% %    }
% %  \resumeSubHeadingListEnd
% 
% 
% %-------------------------------------------
\end{document}
